\chapter{QUẢN LÝ TÍCH HỢP}

Quản lý tích hợp là quá trình hợp nhất các kế hoạch, quyết định và dòng công việc riêng lẻ thành một chỉnh thể vận hành thống nhất. Đối với FundChain, đây là một nội dung đặc biệt quan trọng vì dự án không chỉ có nhiều kế hoạch quản trị song song như phạm vi, lịch biểu, chi phí, chất lượng, rủi ro và truyền thông, mà còn có nhiều dependency kỹ thuật giữa smart contract, backend, frontend, deployment và tài liệu. Nếu các mảnh này không được tích hợp bằng một cơ chế quản trị rõ ràng, dự án sẽ dễ rơi vào tình trạng mỗi nhóm tối ưu phần việc của mình nhưng kết quả chung lại thiếu đồng bộ.

Trong phạm vi môn Quản trị dự án, chương này không chỉ giải thích quản lý tích hợp theo lý thuyết PMBOK, mà còn cụ thể hóa cách FundChain phát triển Project Charter, Project Management Plan, điều phối công việc, kiểm soát thay đổi tích hợp, quản lý cấu hình và kết nối luồng kỹ thuật thành một hệ thống có thể bàn giao. Đây là chương đóng vai trò ``nối'' tất cả các baseline đã xây dựng ở các chương trước.

\iffalse
\section{Tổng quan về quản lý tích hợp}

\subsection{Vai trò của quản lý tích hợp}

Quản lý tích hợp giúp PM nhìn dự án như một hệ thống hoàn chỉnh thay vì tập hợp rời rạc của nhiều hoạt động. Nhờ đó, một thay đổi không còn được đánh giá đơn lẻ mà được xem xét đồng thời dưới góc độ phạm vi, lịch biểu, chi phí, chất lượng, rủi ro, tài nguyên và truyền thông.

Với FundChain, quản lý tích hợp đặc biệt cần thiết ở các tình huống như:

\begin{itemize}
  \item thay đổi contract dẫn đến đổi ABI, cập nhật backend, frontend, test và tài liệu;
  \item rủi ro tích hợp làm lệch cả lịch biểu và chi phí;
  \item quyết định cắt hoặc thêm tính năng đòi hỏi cập nhật nhiều baseline cùng lúc;
  \item deployment thay đổi address hoặc config, kéo theo update toàn bộ môi trường trình diễn.
\end{itemize}

\subsection{Các quy trình tích hợp chính}

Theo hướng tiếp cận quản trị dự án, các quy trình tích hợp chính gồm:

\begin{enumerate}
  \item Develop Project Charter.
  \item Develop Project Management Plan.
  \item Direct and Manage Project Work.
  \item Manage Project Knowledge.
  \item Monitor and Control Project Work.
  \item Perform Integrated Change Control.
  \item Close Project or Phase.
\end{enumerate}

\figureplaceholder{ch11-integration-management-process.png}{Sơ đồ các quy trình quản lý tích hợp: Project Charter, Project Management Plan, Direct and Manage Project Work, Manage Project Knowledge, Monitor and Control Project Work, Integrated Change Control và Close Project/Phase; thể hiện PM ở vai trò trung tâm điều phối.}{Quy trình quản lý tích hợp của dự án}{fig:integration-process}

\section{Triển khai quản lý tích hợp trong dự án}

\subsection{Phát triển Project Charter}

Project Charter là tài liệu khởi đầu giúp dự án xác định rõ lý do tồn tại, mục tiêu, ranh giới và quyền hạn quản trị. Đối với FundChain, Charter cần bao gồm:

\begin{itemize}
  \item nhu cầu dự án và bối cảnh môn học;
  \item mục tiêu SMART;
  \item phạm vi cấp cao và deliverable chính;
  \item milestone và ngân sách cấp cao;
  \item giả định, ràng buộc và rủi ro ban đầu;
  \item stakeholder chính;
  \item quyền hạn điều phối của PM trong phạm vi môn học.
\end{itemize}

PM có thể hỗ trợ soạn thảo Charter, nhưng sponsor/giảng viên là người ban hành hoặc phê duyệt để chính thức cho phép dự án tồn tại và trao quyền cho PM. Charter không phải tài liệu do PM tự phê duyệt.

\subsection{Phát triển Project Management Plan}

Project Management Plan là nơi tích hợp toàn bộ các kế hoạch thành một baseline quản trị thống nhất. Đối với FundChain, kế hoạch này bao gồm tối thiểu:

\begin{itemize}
  \item Scope Baseline;
  \item Schedule Baseline;
  \item Cost Baseline;
  \item Quality Management Plan;
  \item Resource Management Plan;
  \item Communications Management Plan;
  \item Risk Management Plan;
  \item Procurement Management Plan;
  \item Change và Configuration Management Plan.
\end{itemize}

\subsection{Định hướng và quản trị công việc}

Direct and Manage Project Work là quá trình biến kế hoạch thành hành động. Trong dự án này, hoạt động đó được thực hiện thông qua:

\begin{itemize}
  \item sprint planning và work authorization;
  \item phân công task theo owner và priority;
  \item cập nhật tiến độ, issue log và decision log;
  \item rà soát deliverable theo quality gate;
  \item phối hợp liên nhóm khi có dependency hoặc blocker.
\end{itemize}

\subsection{Quản lý tri thức dự án}

Tri thức dự án cần được chuyển từ cá nhân thành tài sản chung. FundChain quản lý tri thức qua:

\begin{itemize}
  \item README, tài liệu kỹ thuật và runbook;
  \item decision record hoặc biên bản cho thay đổi quan trọng;
  \item code review, demo và retrospective;
  \item lessons learned register được cập nhật xuyên suốt dự án.
\end{itemize}

\fi
\section{Theo dõi, kiểm soát và thay đổi tích hợp}

\subsection{Theo dõi và kiểm soát công việc}

PM theo dõi dự án bằng cách đối chiếu actual với baseline ở các góc độ:

\begin{itemize}
  \item scope: deliverable planned so với delivered;
  \item schedule: milestone, SPI và độ lệch activity;
  \item cost: PV, EV, AC, CPI và CV;
  \item quality: quality gate, defect trend và UAT result;
  \item risk: trigger, residual risk và incident;
  \item resource: overload, key-person risk và availability.
\end{itemize}

\begin{table}[H]
\centering
\caption{Dashboard kiểm soát dự án tại data date 31/05/2026}
\label{tab:project-control-dashboard}
\small
\begin{tabularx}{\textwidth}{>{\bfseries}p{2.1cm}p{2.4cm}p{1.4cm}X>{\raggedright\arraybackslash}p{3cm}}
\toprule
Nhóm kiểm soát & Chỉ số/bằng chứng & Trạng thái & Diễn giải thực nghiệm & Hành động tiếp theo \\
\midrule
Scope & Major change = 0 & Xanh & Yêu cầu KYC được ghi backlog, chưa đưa vào release ba tháng & Giữ Scope Baseline, chỉ nhận change qua CCB \\
Schedule & SPI = 0,933 & Vàng & Integration gate chậm do sửa contract và đồng bộ ABI/address & Ưu tiên activity đường găng, dùng crash option đã phân tích \\
Cost & CPI = 0,955 & Vàng nhẹ & Effort tăng vì rework nhưng chưa vượt ngưỡng đỏ & Theo dõi AC tuần kế và không dùng reserve nếu chưa có trigger \\
Quality & 0 Critical/High & Vàng & Còn lỗi Medium về validation, tài liệu và tích hợp & Regression test, UAT checklist và defect triage hằng ngày \\
Risk & R03/R09 tăng & Vàng & Rủi ro redeploy và defect lọt sang UAT cao hơn baseline & Kích hoạt response plan, cập nhật risk owner \\
Procurement & RPC/IPFS trong quota & Xanh & Free tier đủ demo, chưa có outage nghiêm trọng & Duy trì fallback endpoint và dữ liệu demo dự phòng \\
\bottomrule
\end{tabularx}
\end{table}

Dashboard trên là điểm nối giữa các công cụ kiểm soát rời rạc như EVM, Risk Register, Quality Gate và Change Log. Khi một dòng chuyển sang đỏ, PM không chỉ báo cáo trạng thái mà phải xác định quyết định quản trị tương ứng: chấp nhận dùng reserve, thay đổi lịch, giảm phạm vi, tăng kiểm thử hoặc đưa vấn đề lên CCB.

\subsection{Kiểm soát thay đổi tích hợp}

Integrated Change Control là nội dung cốt lõi nhất của chương này. Một change request chỉ được phê duyệt khi tác động của nó được phân tích đồng thời trên nhiều chiều. Bảng~\ref{tab:change-impact-dimensions} mô tả các chiều tác động cần xem xét.

\begin{table}[H]
\centering
\caption{Các chiều tác động cần xem xét trong Integrated Change Control}
\label{tab:change-impact-dimensions}
\begin{tabularx}{\textwidth}{>{\bfseries}p{4cm}X}
\toprule
Chiều tác động & Câu hỏi đánh giá \\
\midrule
Giá trị nghiệp vụ & Thay đổi có thực sự cần thiết với mục tiêu dự án không? \\
Phạm vi & Có làm thay đổi deliverable hoặc WBS đã baseline không? \\
Lịch biểu & Có ảnh hưởng activity đường găng, milestone hoặc buffer không? \\
Chi phí & Có làm tăng effort, reserve hoặc Project Budget không? \\
Chất lượng & Có đòi hỏi thêm test, review hoặc quality gate mới không? \\
Rủi ro & Có tạo rủi ro mới hoặc làm tăng residual risk không? \\
Tích hợp kỹ thuật & Có ảnh hưởng contract storage, ABI, API, config hoặc deployment không? \\
Tài liệu và truyền thông & Có yêu cầu cập nhật tài liệu, hướng dẫn hoặc status report không? \\
\bottomrule
\end{tabularx}
\end{table}

\subsection{Configuration Management}

Configuration Management là một phần thiết yếu của quản lý tích hợp, đặc biệt với dự án có nhiều artifact kỹ thuật như FundChain. Dự án quản lý cấu hình theo các nguyên tắc sau:

\begin{itemize}
  \item source code được version bằng Git và gắn release tag;
  \item dependency được khóa phiên bản có kiểm soát;
  \item ABI được quản lý theo contract version;
  \item network, chain ID, contract address và deployment record phải được lưu có hệ thống;
  \item `.env.example` được duy trì, nhưng secret thật không được commit;
  \item tài liệu triển khai phải cập nhật cùng phiên bản cấu hình tương ứng.
\end{itemize}

\begin{table}[H]
\centering
\caption{Các đối tượng cấu hình trọng yếu của dự án}
\label{tab:configuration-items}
\begin{tabularx}{\textwidth}{>{\bfseries}p{4cm}X}
\toprule
Đối tượng cấu hình & Yêu cầu kiểm soát \\
\midrule
Source code & Version, branch, release tag và review history \\
ABI và contract artifact & Đồng bộ với contract version và frontend/backend integration \\
Deployment record & Network, address, verify link, thời điểm deploy và owner \\
Environment configuration & Template rõ ràng, không chứa secret thật \\
Docker/config runtime & Theo dõi image, service và biến môi trường liên quan \\
Tài liệu hướng dẫn & Cập nhật đồng bộ với release và config hiện hành \\
\bottomrule
\end{tabularx}
\end{table}

\iffalse
\subsection{Chuỗi truy vết cấu hình blockchain}

Với smart contract, một release không được nhận diện chỉ bằng tên branch hoặc tệp ABI. Baseline cấu hình phải truy vết được toàn bộ chuỗi:

\begin{center}
Contract source $\rightarrow$ compiler/settings $\rightarrow$ bytecode $\rightarrow$ ABI $\rightarrow$ deployment transaction $\rightarrow$ address $\rightarrow$ application config $\rightarrow$ integration test $\rightarrow$ release tag.
\end{center}

Contract sử dụng semantic version ở mức release. Thay đổi làm khác ABI hoặc hành vi nghiệp vụ phải tăng version và tạo deployment record mới; không ghi đè address của release cũ. ABI được sinh từ đúng commit đã deploy, lưu cùng artifact và checksum, sau đó được frontend/backend lấy từ một nguồn versioned duy nhất. CI từ chối build nếu checksum ABI, chain ID hoặc address manifest không khớp với release manifest.

\begin{landscape}
\begin{longtable}{p{3.4cm}p{6cm}p{5cm}p{3cm}p{3.2cm}}
\caption{Configuration Status Accounting cho release blockchain}\label{tab:blockchain-csa}\\
\toprule
\textbf{Configuration item} & \textbf{Thuộc tính phải ghi} & \textbf{Kiểm soát/quality gate} & \textbf{Owner} & \textbf{Bằng chứng} \\
\midrule
\endfirsthead
\toprule
\textbf{Configuration item} & \textbf{Thuộc tính phải ghi} & \textbf{Kiểm soát/quality gate} & \textbf{Owner} & \textbf{Bằng chứng} \\
\midrule
\endhead
Contract source & Git commit, tag, compiler, optimizer, constructor args & Review chéo; test và static analysis đạt ngưỡng & BC/TL & Commit, PR, test report \\
ABI/artifact & Contract version, checksum, generated-at commit & Không sửa tay; CI so checksum với artifact & BC & Artifact bundle \\
Deployment record & Network, chain ID, address, deployer, tx hash, block, thời điểm & Receipt thành công; source verified; smoke test đạt & DevOps/TL & Manifest, explorer link \\
Application config & ABI version, address, RPC alias, feature flags & Frontend/backend cùng manifest; integration test đạt & BE/FE & Build log, test report \\
Secret/credential & Tên secret, mục đích, custodian, scope, ngày rotate & Không xuất hiện trong source/log/artifact; least privilege & DevOps/PM & Secret-scan log, access record \\
Release package & Tag, image digest, manifest, runbook, known issues & Approval của TL, QA và PM trước UAT & PM/QA & Release checklist \\
\bottomrule
\end{longtable}
\end{landscape}

\fi
\subsection{Bảo vệ private key và secret trong CI/CD}

Private key của deployer/relayer không được lưu trong repository, tệp \texttt{.env}, image, cache, log hoặc artifact. Pipeline chỉ tham chiếu tên secret từ secret store của nền tảng CI; secret được cấp theo environment, giới hạn quyền đọc và chỉ khả dụng ở job deploy được bảo vệ. Pull request từ fork và job test thông thường không được cấp production-like secret. Log phải mask credential và không in raw transaction signing payload có dữ liệu nhạy cảm.

Tài khoản deployer và relayer tách biệt; ví chỉ giữ lượng Sepolia ETH cần thiết cho thử nghiệm. Thao tác deploy cần reviewer thứ hai phê duyệt, environment protection và ghi lại người thực hiện. Khi nghi ngờ lộ khóa, nhóm dừng job liên quan, thu hồi quyền, rotate secret, chuyển tài sản thử nghiệm nếu cần, redeploy bằng khóa mới, cập nhật manifest và chạy lại regression. Việc chỉ xóa secret khỏi commit gần nhất không được xem là khắc phục đầy đủ vì giá trị có thể còn trong lịch sử Git hoặc log.

\subsection{Quy trình thay đổi và redeploy contract}

Mọi thay đổi sau khi đã deploy phải đi qua các bước: lập Change Request; phân tích tác động storage/interface; review contract và migration; chạy unit/security/regression test; phê duyệt; deploy địa chỉ mới; verify source; cập nhật ABI/address manifest; cập nhật frontend/backend/listener; smoke test; thông báo stakeholder; và lưu release cũ ở trạng thái deprecated. Nếu release mới lỗi, fallback là trỏ ứng dụng về manifest đã phê duyệt gần nhất khi nghiệp vụ cho phép; không giả định có thể ``rollback'' trạng thái on-chain như cơ sở dữ liệu truyền thống.

Checklist trước UAT tối thiểu xác nhận đúng chain ID Sepolia, address và ABI checksum; source đã verify; deployer/relayer đúng quyền; không có secret trong bundle; event listener đọc đúng deployment block; các luồng tài chính trọng yếu vượt smoke test; explorer link và transaction hash đã được ghi vào hồ sơ bàn giao.

\iffalse
\section{Tích hợp kỹ thuật hệ thống}

\subsection{Luồng tích hợp tổng thể}

Ở góc nhìn tích hợp kỹ thuật, FundChain hoạt động theo luồng tổng quát:

\begin{enumerate}
  \item frontend tạo transaction trực tiếp hoặc gửi intent;
  \item backend xử lý phần IPFS, relayer hoặc orchestration ngoài chuỗi khi cần;
  \item contract nền tảng cập nhật trạng thái on-chain;
  \item listener nhận event và đồng bộ trạng thái;
  \item dữ liệu bổ trợ được lưu và phát thông báo;
  \item frontend re-fetch dữ liệu và hiển thị trạng thái cuối cho người dùng.
\end{enumerate}

\figureplaceholder{ch11-system-integration-sequence.png}{Sequence diagram mô tả hai luồng: direct transaction và gasless transaction. Các thành phần gồm User/Frontend, Wallet, Backend/Relayer, IPFS, Smart Contract và Event Listener/Notification.}{Luồng tích hợp tổng thể của hệ thống FundChain}{fig:system-integration-sequence}

\fi
\subsection{Integration gate}

Để hạn chế lỗi dây chuyền, dự án sử dụng integration gate trước khi chuyển sang UAT:

\begin{itemize}
  \item contract core đã ổn định và vượt test cần thiết;
  \item backend và frontend dùng đúng ABI/address/config;
  \item luồng E2E trọng yếu chạy được trong môi trường thử nghiệm;
  \item tài liệu deploy và known issues được cập nhật;
  \item defect Critical/High liên quan đến tích hợp đã được xử lý.
\end{itemize}

\section{Đánh giá tài chính bằng NPV}

\subsection{Giả định phân tích}

Để minh họa góc nhìn đầu tư ở mức quản trị, dự án sử dụng các giả định sau:

\begin{itemize}
  \item initial investment quy đổi bằng Project Budget: 64.000.000 VNĐ;
  \item lợi ích quy đổi giả định: 30.000.000 VNĐ/năm trong 3 năm;
  \item discount rate giả định: 12\%/năm;
  \item không tính terminal value.
\end{itemize}

Các biến trên thuộc nhóm \textit{simulated assumptions}, không phải doanh thu hoặc dòng tiền đã quan sát. Báo cáo phân biệt ba loại dữ liệu: \textit{planned} là baseline được phê duyệt; \textit{simulated} là số liệu giả định dùng để thực hành phương pháp; \textit{actual} là dữ liệu có chứng từ hoặc log phát sinh trong quá trình thực hiện. EVM và NPV không được trình bày như kết quả thương mại nếu chưa có actual cost và actual benefit tương ứng.

\subsection{Kết quả NPV}

Theo giả định trên:

\begin{center}
$NPV = -64.000.000 + \frac{30.000.000}{1{,}12} + \frac{30.000.000}{1{,}12^2} + \frac{30.000.000}{1{,}12^3}$
\end{center}

Kết quả xấp xỉ:

\begin{center}
\textbf{NPV $\approx$ 8.055.000 VNĐ}
\end{center}

NPV dương cho thấy, trong kịch bản giả định, phương án đầu tư có thể chấp nhận được về mặt tài chính. Tuy nhiên, cần nhấn mạnh rằng đây không phải dự báo thương mại thực tế mà chỉ là mô phỏng học thuật để hỗ trợ đánh giá tích hợp giữa đầu tư, chi phí và lợi ích.

\begin{table}[H]
\centering
\caption{Phân tích nhạy cảm định lượng cho NPV}
\label{tab:npv-sensitivity}
\begin{tabularx}{\textwidth}{>{\bfseries}p{3cm}p{3cm}p{3.2cm}X}
\toprule
Lợi ích/năm & Discount rate & NPV & Kết luận kịch bản \\
\midrule
20 triệu VNĐ & 12\% & $-15{,}963$ triệu & Âm; không đạt điều kiện tài chính giả định \\
30 triệu VNĐ & 12\% & $8{,}055$ triệu & Dương nhẹ; kịch bản cơ sở \\
40 triệu VNĐ & 12\% & $32{,}073$ triệu & Dương; có biên an toàn tốt hơn \\
30 triệu VNĐ & 15\% & $4{,}497$ triệu & Vẫn dương nhưng nhạy với discount rate \\
30 triệu VNĐ & 10\% & $10{,}606$ triệu & Dương; thuận lợi hơn kịch bản cơ sở \\
\bottomrule
\end{tabularx}
\end{table}

Tại discount rate 12\%, lợi ích hòa vốn hằng năm là $64/2{,}401831\approx26{,}646$ triệu VNĐ. Đây là ngưỡng của mô hình ba năm đã giả định, không phải forecast doanh thu. Nếu chuyển sang production, mô hình phải bổ sung audit, pháp lý, gas Mainnet, vận hành/SLA, incident response, bảo hiểm và dữ liệu adoption thực tế trước khi tái tính NPV.

\iffalse
\section{Công cụ hỗ trợ và bài học tích hợp}

\subsection{Công cụ hỗ trợ}

Các công cụ hỗ trợ tích hợp của dự án gồm:

\begin{itemize}
  \item Git/GitHub cho source, review và tag release;
  \item bộ công cụ test, build và coverage cho kiểm soát chất lượng;
  \item Docker hoặc môi trường đóng gói tương đương cho triển khai thử nghiệm;
  \item các dịch vụ RPC, IPFS, database, queue và hosting;
  \item GitHub Issues/Projects, Google Drive, chat nhóm và sơ đồ thiết kế.
\end{itemize}

\subsection{Thách thức và bài học kinh nghiệm}

Các thách thức tích hợp chính của FundChain gồm:

\begin{itemize}
  \item tài liệu dễ bị chậm cập nhật so với source;
  \item breaking change ở contract có tác động dây chuyền lớn;
  \item cần phân biệt rõ on-chain là source of truth tài chính còn off-chain phục vụ UX và truy vấn;
  \item dependency bên ngoài có thể gây gián đoạn dù logic nội bộ đúng;
  \item việc giảm gas không được đánh đổi bằng giảm an toàn tài chính hoặc bỏ test.
\end{itemize}

\fi
\section{Kết luận chương}

Chương XI đã tích hợp các kế hoạch riêng lẻ của dự án thành một khung điều phối thống nhất, từ Charter, Project Management Plan, Direct and Manage Work, Configuration Management, Integrated Change Control đến tích hợp kỹ thuật hệ thống và phân tích NPV minh họa. Nhờ đó, quản lý tích hợp được thể hiện đúng vai trò là ``lớp điều phối trung tâm'' của toàn bộ dự án.

Đối với FundChain, quản lý tích hợp là điều kiện để mọi thay đổi đều được xem xét đầy đủ trước khi triển khai và để các thành phần kỹ thuật vận hành như một chỉnh thể. Đây cũng là nền tảng trực tiếp cho việc kết thúc dự án một cách có kiểm soát ở chương cuối.
